% Lösungen Einheit 3 von 8 – Pfadregeln bei unabhängigen Stufen (zusammenbau v0.1)
\einheitenkopf{Einheit 3 von 8 · Pfadregeln bei unabhängigen Stufen}

\erg{17}{a) p bleibt gleich \quad b) Äste R $0{,}3$, B $0{,}7$ auf beiden Stufen; $P = 0{,}3 \cdot 0{,}7 = 0{,}21$ \quad c) $0{,}6^3 \cdot 0{,}4^2 \approx 0{,}0346$ \quad d) $0{,}35^4 \approx 0{,}0150$ \quad e) $\frac{9}{64} + \frac{25}{64} = \frac{17}{32}$ \quad f) P(A) $= 0{,}42$; P(B) $= 1 - 0{,}49 = 0{,}51$; Gegenereignis: beide Male Niete, $P = 0{,}49$ \quad g) $1 - 0{,}85^6 \approx 0{,}623$ \quad h) $\frac{1}{6} + \frac{5}{36} + \frac{25}{216} = \frac{91}{216}$ \quad i) P(A) $= \frac{61}{216} \approx 0{,}282$; P(B) $= \frac{305}{1\,296} \approx 0{,}235$; A ist im Vorteil \quad j) eine Runde: $\frac{1}{2} + \frac{1}{2} \cdot \frac{1}{2} = \frac{3}{4}$; zwei Runden: $\left(\frac{3}{4}\right)^2 = \frac{9}{16}$}

17b)

\baumzwei{R/0{,}3,B/0{,}7}{R/0{,}3,B/0{,}7}{R/0{,}3,B/0{,}7}


17f)

\baumzwei{G/0{,}3,N/0{,}7}{G/0{,}3,N/0{,}7}{G/0{,}3,N/0{,}7}

\erg{18}{a) $\frac{1}{6}$: (X; rot), (Y; blau), (Y; weiß)}
\erg{19}{a) A: $0{,}40 : 0{,}72 \approx 0{,}56$\,€; B: $0{,}30 : 0{,}525 \approx 0{,}57$\,€; A ist günstiger}
\erg{20}{a) Ole vervielfacht statt zu multiplizieren; Pfadregel: $0{,}3^3 = 0{,}027$ \quad b) Nach dem Zurücklegen ist die Urne wieder wie am Anfang; jede Stufe hat dieselben Anteile. \quad c) $1 - 0{,}05^2 = 0{,}9975$ \quad d) $0{,}6 \cdot 0{,}4 + 0{,}4 \cdot 0{,}6 = 0{,}48$}
